\begin{abstract}
This paper presents a lightweight browser-based 3D engine designed to support the development and execution of interactive educational scenarios without requiring native runtime installation. The engine is needed because existing approaches to delivering 3D learning content through the browser force developers to choose between heavyweight exported runtimes with long loading times and low-level graphics libraries that demand specialized rendering expertise. The proposed solution provides a high-level component model over a hidden rendering backend, enabling educational platform developers to build interactive 3D scenarios with reduced engineering effort while maintaining runtime efficiency on consumer hardware. The paper describes the key architectural choices behind the engine and introduces a set of targeted optimization strategies addressing asset memory consumption, rendering initialization latency, and per-frame computational overhead. The proposed strategies are intended to be assessed on three benchmark scenes of increasing complexity -- a procedural geometry grid, a high-polygon model, and an interactive Solar System educational scenario -- using frame time, loading latency, and heap memory as the main evaluation metrics. This evaluation is intended to determine the practical suitability of the proposed approach for browser-based educational 3D applications.

\keywords{WebGL \and game engine architecture \and interactive learning \and 3D visualization \and performance optimization \and technology-enhanced learning}

\end{abstract}